\documentclass[11pt,a4paper]{article}
\usepackage[top=1.00in, bottom=1.0in, left=1.1in, right=1.1in]{geometry}
\usepackage{hyperref}
\usepackage[style=authoryear, backend=biber, natbib = true, style=numeric-comp, sortcites, hyperref, sorting=none, giveninits=true, doi=false,isbn=false,url=false,eprint=false, date=year, maxbibnames = 2, backend=bibtex]{biblatex}
\AtEveryBibitem{\clearfield{pages}}
\AtEveryBibitem{\clearlist{language}}
\AtEveryBibitem{%
  \clearfield{volume}%
  \clearfield{number}}
\AtEveryBibitem{\clearfield{title}}
\renewbibmacro{in:}{}

\addbibresource{/home/victor/projects/synchrony/docs/pnas/synchrony.bib} 

\title{Reply to Michał Bogdziewicz \emph{et al.}: Research on the summer solstice needs falsifiable predictions} % Title! Research on the summer solstice needs falsifiable predictions

\author{Victor, Lizzie} % emw4July -- make this look a little more professional before submitting
\date{August 2024}

\begin{document}

\maketitle

In \citep{VanderMeersch2025} we showed that summer solstice appears as an optimal time to transition between major physiological growth and reproductive phases. Given a trade-off between predicting total and remaining thermal heat, summer solstice was optimal on average for Europe and North America, but with local variation. Bogdziewicz \emph{et al.} interpret our results as invalidating the potential role of solstice in mast seeding \citep{Bogdziewiczresponse}, but we suggest they have mis-interpreted our work. We explicitly considered whether solstice may control the timing of mast seeding, and our results do not invalidate this hypothesis. In \citep{VanderMeersch2025} we do, however, consider the potential role of other drivers and raise questions about how to differentiate between them given strong covariation. 
% Our results do show solstice is optimal, but also highlight some amount of variation across space, and time.

% Zohner is really 1-2 weeks AFTER and Journe is BEFORE
Bogdziewicz \emph{et al.} argue that summer solstice is a better cue for masting than temperature-related cues because it allows for greater synchrony across individuals \citep{Bogdziewiczresponse}, whereas we highlight spatial variation in a trade-off based on temperature accumulation. We do not disagree that we show important variability across space. We question, however, their conclusions because: (i) in many masting trees, the strength of synchrony among individuals decays with distance, inconsistent with cuing directly to summer solstice \citep{Liebhold2004}, and (ii) the strong covariation between summer solstice and balance in annual heat accumulation makes it difficult to reject either as the primary cue---a major message of our paper. Showing a masting pattern consistent with cuing to solstice is not the same as demonstrating trees cue to solstice.

In \citep{VanderMeersch2025} we discuss the challenge in differentiating between cues. Current work on whether the summer solstice acts as ``\emph{celestial starting gun}'' %that starts a ``\emph{temperature-sensing window}"
has focused on a statistical correlation between an annual metric of seed production and daily temperature to show a shift in the correlation may occur near the summer solstice on average, but this timing similarly varies across space (\citep{Journe2024}, Fig. 1 in \citep{Bogdziewiczresponse}). These results echo ours \citep{VanderMeersch2025} and thus we are confused as to where Bogdziewicz \emph{et al.} see a discrepancy between these findings and ours. Identifying what level of variation invalidates the solstice-as-cue hypothesis seems an important question, especially as results from summer solstice research for both masting \citep{Journe2024,Bogdziewiczresponse} and growth \citep{Zohner2023} show the statistically most likely date before or after the solstice by one or more weeks. % (e.g., Fig. 1 in submitted response where the `window of opportunity' appears to begin around XX). % Cut words in () here at end?

% Maybe we cut the below paragraph for space? If we want we could sneak some of the second half of it into the second paragraph.
% Developing testable predictions for variation across space could help to understand masting.
% Based on Bogdziewicz \emph{et al.} plants across the northern hemisphere would be predicted to mast seed in synchrony \citep{Bogdziewiczresponse}, but this is not common (findCITES).
% While we agree that many economy of scale hypotheses for masting require some level of synchrony to lead to increased fitness, most do not predict extremely large-scale synchrony \citep{Davies2024} and may thus predict cues other than solstice. Even in the case of extremely widespread synchrony, whether the solstice alone would be an ideal cue for masting is not clear from current work. 

Developing falsifiable predictions to test the role of the solstice---and other cues---in driving transitions between major physiological phases is important for fundamental science and forecasting. Identifying the spatial scale of synchrony that maximizes fitness based on different economy of scale hypotheses for masting \citep{Davies2024} could rule out some cues. We hope more efforts that build from eco-evolutionary and climatological perspectives \citep{VanderMeersch2025} can provide clearer predictions and stronger evidence, especially as bounded systems (e.g., the growing season, \citep{Colwell2000, Morales2005}) and natural climatological trends often lead to interesting patterns with multiple competing explanations that make very different predictions \citep{Gao2024,decsens}. % for broken stick: https://www.cell.com/trends/ecology-evolution/abstract/S0169-5347(99)01767-X?large_figure=true and/or Morales2005
% Our goal in (citeus) was to encourage this discussion by to consider the link between the solstice and these processes through a broader lens---both from a climatological and an eco-evolutionary perspectives. Our work did not aim at refuting any specific hypothesis. Instead, we believe it emphasizes that robust science depends on considering multiple mechanisms and alternative hypotheses---always with a particular caution about any causal links we believe we are testing. 

\paragraph{References}
\printbibliography[heading=none]

\end{document}
